\section{Part 2：Neural 和 LLM scaling behaviors}


\singleslide{slides-images/slide-012.jpg}{Part 2 的三个问题：data→performance、data vs model allocation、hyperparameters→performance。}{12}

Slide 12 是全讲实验地图。第一问固定模型看数据规模，第二问在 compute budget 下权衡更多 tokens 与更多 parameters，第三问研究 architecture、optimizer、batch 等超参数能否从小规模迁移。三类曲线共同把“大模型设计”变成受约束预测问题。

\begin{knowledgebox}{讲义补充：源 Slide 12 的三个问题}
第一，数据量如何影响 loss；第二，给定 compute 时应该增加数据还是模型参数；第三，架构、优化器、batch size、learning rate 等超参数如何随规模变化。这三类问题对应 LLM 训练前最昂贵的决策。
\end{knowledgebox}

\begin{figure}[H]
\centering
\includegraphics[width=0.86\textwidth]{slide-013.jpg}
\caption{Slide 13：Scaling laws 在很多因素和非标准设置中都呈 power-law 关系。}
\figsource{CS336 Spring 2026 Lecture 9 官方幻灯片。}
\end{figure}

\begin{knowledgebox}{读图：Slide 13 强调“many factors”}
Power law 不只出现在数据量上，也出现在模型参数、计算量、某些架构选择和任务指标上。它的工程吸引力来自同一种拟合工具可以服务多个设计问题。但这也带来风险：不是所有指标、所有下游任务都同样平滑。
\end{knowledgebox}

\subsection{data vs performance}

本节把最基础的自变量固定为训练数据量，观察验证 error 或 cross-entropy loss 如何变化。核心不是看到一条直线就宣布规律成立，而是明确坐标变换、不可约损失、有效数据定义和拟合区间，再检查同一 exponent 是否跨模型规模、domain 与训练 recipe 稳定。

\begin{figure}[H]
\centering
\includegraphics[width=0.86\textwidth]{slide-014.jpg}
\caption{Slide 14：data scaling law 把 dataset size \(n\) 映射到 error，期望曲线单调且类似 logistic/power-law。}
\figsource{CS336 Spring 2026 Lecture 9 官方幻灯片。}
\end{figure}

\begin{knowledgebox}{读图：Slide 14 的数据 scaling law}
横轴是数据集大小 \(n\)，纵轴是 error 或 loss。合理曲线应该单调下降，但不可能无限下降到负数；在有限区间内，power law 常能很好拟合。曲线的 slope 告诉我们继续加数据的边际收益。
\end{knowledgebox}

\begin{figure}[H]
\centering
\includegraphics[width=0.86\textwidth]{slide-015.jpg}
\caption{Slide 15：语言模型中 test loss 和 dataset size 在 log-log 图上近似线性，即 scale-free/power law。}
\figsource{CS336 Spring 2026 Lecture 9 官方幻灯片。}
\end{figure}

\begin{importantbox}{读图：Slide 15 的 log-log 直线}
如果
\[
L(n)-L_\infty = A n^{-\alpha},
\]
两边取对数得到
\[
\log(L(n)-L_\infty)=\log A-\alpha\log n.
\]
其中 \(L(n)\) 是数据量为 \(n\) 时的 loss，\(L_\infty\) 是不可约或远端 offset，\(\alpha\) 是数据 scaling exponent。log-log 图上的斜率就是 \(-\alpha\)。
\end{importantbox}


\singleslide{slides-images/slide-016.jpg}{Data scaling 的概念基础：从均值估计例子理解 estimation error 为什么多项式衰减。}{16}

Slide 16 先问“为什么是 power law”，再用最简单的 mean estimation 建立直觉。独立样本平均的标准误差随 $n^{-1/2}$ 衰减，在 log-log 图上就是直线。神经网络指数未必等于 $1/2$，但 estimation error 的多项式收敛解释了直线形式为何自然出现。

\begin{warningbox}{讲义补充：源 Slide 16 的解释不是完整理论}
“估计误差多项式衰减”能解释为什么 power law 合理，但不能完全解释神经网络 exponent 的数值、不同数据分布的 offset、下游任务的不平滑性。Scaling laws 在工程上有用，不代表我们已完全理解其机制。
\end{warningbox}

\begin{figure}[H]
\centering
\includegraphics[width=0.86\textwidth]{slide-017.jpg}
\caption{Slide 17：mean estimation toy example，误差 \(\mathbb{E}[(\hat{\mu}-\mu)^2]=\sigma^2/n\) 是最简单 scaling law。}
\figsource{CS336 Spring 2026 Lecture 9 官方幻灯片。}
\end{figure}

\begin{importantbox}{读公式：Slide 17 的 mean estimation}
输入 \(x_1,\dots,x_n\sim\mathcal{N}(\mu,\sigma^2)\)，估计量 \(\hat{\mu}=\frac{1}{n}\sum_i x_i\)。均方误差为
\[
\mathbb{E}\left[(\hat{\mu}-\mu)^2\right]=\frac{\sigma^2}{n}.
\]
取对数后是 \(\log \text{Error}=-\log n + 2\log\sigma\)，斜率为 \(-1\)。这说明 power law 可以从很基础的统计估计中出现。
\end{importantbox}


\singleslide{slides-images/slide-018.jpg}{Scaling exponent 之谜：经典参数模型常预期 $1/n$，机器翻译、语音和语言建模却显示不同斜率。}{18}

这页把“线性”与“斜率”分开。Power law 形式可以普遍存在，但指数编码任务维度、噪声、模型适配与数据分布；真实神经任务远慢于简单参数估计的 $n^{-1}$。因此不能拿一个领域的 exponent 直接预测另一个领域。

\begin{knowledgebox}{讲义补充：源 Slide 18 的 mystery}
若所有任务都像均值估计，log-log 斜率应接近 \(-1\)。但机器翻译、语音、语言模型等神经 scaling 的 slope 通常更平缓。这说明模型、数据分布、任务复杂度和有效维度共同决定 exponent，而不是单纯样本平均。
\end{knowledgebox}


\singleslide{slides-images/slide-019.jpg}{非参数学习直觉：将输入空间分箱后，误差指数随 intrinsic dimension 变成 $n^{-1/d}$。}{19}

Slide 19 用二维函数估计说明灵活模型的样本要覆盖空间。维度 $d$ 越高，局部区域获得足够样本越慢，收敛指数越小；神经网络若利用低维结构，则有效 $d$ 可能远低于输入维度。这个例子不是 Transformer 理论，而是解释不同任务斜率的候选机制。

\begin{knowledgebox}{讲义补充：源 Slide 19 的非参数直觉}
如果目标函数在二维空间上变化，把空间切成更细网格需要更多样本才能估计每个区域。维度越高，样本需求越大，error 下降越慢。这给 scaling exponent 和 intrinsic dimensionality 的联系提供直觉。
\end{knowledgebox}

\begin{figure}[H]
\centering
\includegraphics[width=0.86\textwidth]{slide-020.jpg}
\caption{Slide 20：intrinsic dimensionality theory 认为 exponent \(\alpha\) 可能与数据内在维度相关。}
\figsource{CS336 Spring 2026 Lecture 9 官方幻灯片。}
\end{figure}

\begin{knowledgebox}{读图：Slide 20 的理论路线}
这页给出一种解释：学习速率 \(n^{-\alpha}\) 中的 \(\alpha\) 受数据或函数的 intrinsic dimension 影响。维度越高，可学习结构越复杂，error 随数据下降越慢。该理论有启发性，但对现代 LLM 的完整解释仍不充分。
\end{knowledgebox}

\subsection{data composition、shift 和 repetition}


\singleslide{slides-images/slide-021.jpg}{从 dataset size 扩展到 composition 与 repetition：数据量相同，mixture 和重复率也会改变性能。}{21}

Slide 21 把标量 $D$ 拆开。真实 pretraining data 由多个 domains 组成，各自质量、难度和可用 token 不同；小模型可用于搜索 mixture，重复高质量数据又会带来记忆与边际收益下降。Scaling law 因而也可以预测“哪些 tokens 值得收集和重放”。

\begin{knowledgebox}{讲义补充：源 Slide 21 把数据问题从“多少”扩展到“什么”}
数据量相同，不同 mixture 可能导致不同 loss offset 和下游能力。Scaling law 不只用来决定 token 数，也可以比较 data mixture、domain mix、是否重复有限数据等问题。
\end{knowledgebox}

\begin{figure}[H]
\centering
\includegraphics[width=0.86\textwidth]{slide-022.jpg}
\caption{Slide 22：distribution shift 下，data composition 可能改变 offset，而不是 slope。}
\figsource{CS336 Spring 2026 Lecture 9 官方幻灯片。}
\end{figure}

\begin{importantbox}{读图：Slide 22 的 offset 工程含义}
如果不同数据 mixture 主要改变 offset，那么小规模实验可以帮助选择“整条曲线更低”的数据组合。Slope 相近意味着继续 scale 时排序可能保持；但如果下游任务或分布改变，offset 和 slope 都可能变化。
\end{importantbox}

\begin{figure}[H]
\centering
\includegraphics[width=0.86\textwidth]{slide-023.jpg}
\caption{Slide 23：实践中用 scaling 做 data mixture selection 很难，小规模最佳不一定等于大规模最佳。}
\figsource{CS336 Spring 2026 Lecture 9 官方幻灯片。}
\end{figure}

\begin{warningbox}{读图：Slide 23 的警告}
自然想法是用小模型拟合每种数据 mixture 的 scaling law，再选最优。但数据质量、任务覆盖、去重、污染、下游迁移、模型大小都可能改变排序。因此 mixture selection 的 scaling 需要多尺度验证，不能只拿最小模型的赢家。
\end{warningbox}

\begin{figure}[H]
\centering
\includegraphics[width=0.86\textwidth]{slide-024.jpg}
\caption{Slide 24：data repetition 下，有效数据 \(D'\)、unique tokens \(U_d\)、重复率 \(R_d\) 的关系。}
\figsource{CS336 Spring 2026 Lecture 9 官方幻灯片。}
\end{figure}

\begin{knowledgebox}{读公式：Slide 24 的重复数据账本}
\(U_d\) 表示 unique tokens，\(R_d\) 表示重复次数，\(D'\) 表示有效数据量。重复数据不是完全没用，但边际价值低于新数据。随着重复率上升，有效数据增长会折损，scaling curve 可能偏离只看 token count 的预测。
\end{knowledgebox}

\begin{figure}[H]
\centering
\includegraphics[width=0.86\textwidth]{slide-025.jpg}
\caption{Slide 25：compute unbounded 设置下 scaling laws 可能 break，因为它们只是 lower bounds。}
\figsource{CS336 Spring 2026 Lecture 9 官方幻灯片。}
\end{figure}

\begin{warningbox}{读图：Slide 25 不要盲目外推}
Scaling laws 是已观察训练 recipe 下的经验下界或趋势。若算法、数据清洗、架构、优化器或训练过程改变，你可能做得比旧曲线更好；若超出拟合区间太远，也可能更差。它们是决策工具，不是自然定律。
\end{warningbox}

\begin{figure}[H]
\centering
\includegraphics[width=0.86\textwidth]{slide-026.jpg}
\caption{Slide 26：考虑数据有限性后，data selection 应随 scale 自适应调整。}
\figsource{CS336 Spring 2026 Lecture 9 官方幻灯片。}
\end{figure}

\begin{knowledgebox}{读图：Slide 26 的 adaptive data selection}
如果重复数据价值下降，那么不同规模下最优数据选择可能不同：小模型需要高质量/高密度信号，大模型可能需要更多覆盖和多样性。数据 pipeline 不应是一次性静态配方，而应随目标 compute 和模型规模变化。
\end{knowledgebox}


\singleslide{slides-images/slide-027.jpg}{Data scaling 小结：log-data 与 log-error 的线性、跨域稳定性、均值估计直觉与数据策展应用。}{27}

这页把 Part 2 的结论闭环：经验 power law 可跨域外推，但指数需要测量；理论例子解释形状而非给出完整神经网络定律；工程上曲线可用于评估新增数据、mixture 与 repetition 的价值。下一部分将同一方法用于 architecture 与 optimizer 决策。

\begin{importantbox}{讲义补充：源 Slide 27 的 recap}
数据 scaling 的核心事实是 log-log 线性；理论直觉来自 generalization/sample-complexity；实践用途包括数据收集、数据重复、data mixture 和分布迁移。它最大的限制是：loss 曲线平滑不代表所有下游能力都同样可预测。
\end{importantbox}

\subsection{本章小结}

数据 scaling laws 告诉我们，更多数据通常按 power law 改善 loss，但数据质量、composition、重复和分布迁移会改变 offset 甚至 slope。用于训练决策时，要同时记录 unique tokens、重复率、mixture、评估分布和模型规模。

